\section{Limitations and Future Work}

The present work has four limitations. First, the evaluated implementation requires an experimental Clang P2996 branch, so deployment depends on the compiler revision recorded with the benchmark artifact. Second, the initial release covers wire serialization only: gRPC service stubs are out of scope, and integration with a gRPC runtime is untested, Section~\ref{sec:interop} covering the framing alone. Third, constant evaluation is bounded first by the compiler's step limit rather than by memory: a single message of 300 \texttt{int32} fields compiles under Clang's default of $1{,}048{,}576$ steps, while 301 fields fail. The bound is a flag rather than a wall, since \texttt{-fconstexpr-steps=4194304} compiles a $1{,}000$-field message, but \texttt{ctpb} does not raise the default. Finally, the performance data comes from one host and a limited set of synthetic schemas and one 47-byte \texttt{Person} message. Future work targets the Well-Known Types (\texttt{Timestamp}, \texttt{Duration}), compile-time JSON mapping, broader workload coverage, and asynchronous stub synthesis.

\section{Conclusion}
This paper presented \texttt{ctpb}, a header-only C++26 library that performs Protocol Buffers serialization without generated C++ source. By combining \texttt{constexpr} recursive-descent parsing with P2996 static reflection, \texttt{ctpb} compiles \texttt{.proto} schemas directly into type-safe C++ structures in the compiler front-end. On one host, a 100-field synthetic schema is parsed in $7.239\,\mu\text{s}$, and the tested 47-byte \texttt{Person} workload encodes $5.38\times$ to $5.59\times$ faster than the full and lite \texttt{libprotobuf} variants. Decoding remains near parity while including UTF-8 validation, and the exercised cases pass bidirectional wire-format interoperability tests against \texttt{libprotobuf}. The static-runtime generated comparison is $93.70\times$ larger than the \texttt{ctpb} sample, while the dynamic generated comparison is $1.46\times$ larger. These results show that static reflection can remove generated Protocol Buffers source for the evaluated schemas and toolchain while retaining tested wire-format interoperability.
